\chapter{GPU 索引设计与构建}
\label{ch:architecture}
压缩 GPU 索引既需要可搜索的表示，也需要能在设备显存内完成的构建过程。本章说明 JHQ 的训练、编码和按列表组织的存储 \cite{jhq}，再检查构建后的表示。查询执行与实测构建成本见 \cref{ch:gpu-execution,sec:7-7}。

\section{构建与搜索流程}
\label{sec:4-1}
train 拟合模型，add 将行编码并按列表排序，search 使用完成的索引。缓存命中跳过拟合，但不跳过编码，因此冷构建与缓存搜索的成本不同。

\begin{figure}[!htbp]
\centering
% Audited lifecycle; no proposed storage or distributed subsystem.
\begin{tikzpicture}[
 font=\small\linespread{1}\selectfont,
 box/.style={draw=ThesisNavy!65,rounded corners=2pt,align=center,
   text width=3.65cm,minimum height=1.9cm,inner sep=5pt},
 gpu/.style={box,fill=ThesisNavy!7},
 arr/.style={-{Latex[length=2mm]},draw=ThesisNavy!80,thick}]
\node[box] (host) at (0,0) {Host input vectors\\Training samples\\Database batches};
\node[box] (model) at (4.8,0) {Host model\\QR rotation $\Pi$\\Primary codebook};
\node[box] (cache) at (9.6,0) {Host model cache\\$\Pi$ and codebooks\\IVF centroids};
\node[gpu] (train) at (0,-3.6) {GPU training\\Rotate samples\\Fit IVF centroids\\Fit residual model};
\node[gpu] (add) at (4.8,-3.6) {GPU bulk build\\1 or 2 passes\\Assign / encode\\Reorder / pack\\Floats $\leq768$ MiB};
\node[gpu] (index) at (9.6,-3.6) {GPU-resident index\\Offsets and IDs\\Primary / residual codes\\Corrections};
\node[box] (query) at (0,-7.2) {Host query batch\\Pinned staging\\Transfer to GPU};
\node[gpu] (search) at (4.8,-7.2) {GPU query pipeline\\Rotate; probe; scan\\Refine; final top-$k$\\Workspace / graph};
\node[box] (result) at (9.6,-7.2) {Host query results\\IDs and scores\\Transfer from GPU};
\draw[arr] (host) -- node[left,font=\footnotesize] {samples} (train);
\draw[arr] (model.south west) -- (train.north east);
\draw[arr] (train) -- node[above,font=\footnotesize] {model} (add);
\draw[arr] (host.west) -- ++(-.5,0) |- (4.8,-5.15)
 node[pos=.8,below,font=\footnotesize] {database batches} -- (add.south);
\draw[arr] (model) -- (add);
\draw[arr] (add) -- (index);
\draw[<->,dashed,draw=ThesisNavy!80,thick]
 (train.north east) -- node[pos=.8,above,sloped,font=\footnotesize,fill=white,inner sep=1pt]
 {save / load model} (cache.south west);
\draw[arr] (query) -- (search);
\draw[arr] (search) -- (result);
\draw[arr] (index.south) -- ++(0,-.6) -| (search.north);
\end{tikzpicture}

\caption[构建与查询架构]{实现中的构建与查询边界。主机工作包括抽样、正交矩阵构造和解析主码本生成；评测配置在 GPU 上拟合 IVF 与残差模型。缓存保存训练状态，完整数据库表示则位于设备上。虚线表示模型复用。暂存上限不包含其他临时空间。}
\label{fig:4-architecture}
\end{figure}

主机准备样本、旋转和主码本；GPU 执行旋转，拟合 IVF／残差模型，并将有界 batch 编码为列表顺序。

模型保存旋转、码本和质心；数据库状态保存 offsets、ID、编码和校正项。CUDA Graph 重放已记录的 GPU 操作，并引用复用的查询缓冲。

一次成功的 \texttt{add} 固定物理位置。后续 add 被拒绝，更新需要重建。可变／事务性发布尚未实现。

主机 ID／分数及调用计时遵循 \cref{eq:5-refined-constant,sec:3-3}。


\section{训练、旋转与码本生成}
\label{sec:4-2}
\subsection{变换与抽样}

主机上的 QR 分解将带种子的高斯矩阵分为正交因子和上三角因子，其中正交因子给出方形 $\Pi$。训练样本和数据库向量使用同一变换：
\begin{equation}
 Y_t=X_t\Pi^{\mathsf T},\qquad y=\Pi x,\qquad
 \|\Pi(q-x)\|_2^2=\|q-x\|_2^2
 \quad\text{when }\Pi^{\mathsf T}\Pi=I.
 \label{eq:4-transform}
\end{equation}
维度仍为 $d$。\cref{sec:4-6} 的验证区分存储矩阵的几何性质（K0）与通用矩阵乘法（general matrix multiplication，GEMM）的舍入。

主评测版本（v69） 存储 FP32 矩阵，但在 $d\geq1{,}536$ 时启用 cuBLAS TF32 tensor-core 运算。这影响两个 OpenAI 数据集在训练／add／查询中的旋转和质心 GEMM；实现提供的环境变量开关无法关闭这一模式。因此，FP32 存储并不表示完全采用 FP32 乘法。FP32 变体（v72） 修正了这一开关。敏感性实验复用同一个 TF32 训练缓存，改变 add／查询算术，并非以 FP32 重新训练模型（\cref{sec:7-precision}）。

主码／IVF 训练从数据库中等间隔选取样本，请求的样本数为 $n_t=39n_{\rm list}$：
\[
 n_t=\begin{cases}159{,}744,&\text{Vogue},\\319{,}488,&\text{arXiv/OpenAI},\\1{,}277{,}952,&\text{BGE/Stella}.\end{cases}
\]
这是 Faiss 的样本数警告阈值，而非量化质量保证 \cite{faiss_clustering}。

\subsection{解析主码本}

主码本由样本尺度解析生成，其 Cartesian 等级保留了编码器与查找表使用的坐标结构：
\begin{equation}
 \sigma^2=\frac{1}{n_t d}\sum_{i=0}^{n_t-1}\|x_i\|_2^2,
 \qquad D_s=d/M,\qquad K=2^B.
 \label{eq:4-scale}
\end{equation}
\[
 B=D_s=8\Longrightarrow
 \mathcal C_m=\{-a,+a\}^{8},\quad a=0.674\sigma.
\]
八个符号可放入一个 byte，支持 $O(d)$ 编码与拆分 LUT \cite{jhq}（\cref{sec:5-2}）。无约束 Lloyd 更新可能破坏这一 Cartesian 结构。对固定符号向量 $s\in\{-1,+1\}^d$，平方距离展开为
\begin{equation}
 \|q'-a s\|^2=\|q'\|^2+da^2-2a\langle q',s\rangle.
 \label{eq:4-sign-ranking}
\end{equation}
\[
 a>0\Longrightarrow
 \arg\min_s\|q'-as\|^2=\arg\max_s\langle q',s\rangle.
\]
因此，在 $a>0$ 时，主码排序由查询与符号向量的内积决定，无需数据服从高斯分布。残差仍依赖 $a$；零尺度和 FP32 并列需要另行处理。

\subsection{粗分配中的筛选与精确计算}

IVF 用间隔选取的旋转行初始化质心，并以 $s_\ell=\|\mu_\ell\|^2-2\langle y,\mu_\ell\rangle$ 运行八次迭代。支持的维度先用 int8 分数筛选候选，再在 FP32 中处理不确定情况。下面的引理说明区间筛选保留分配结果所需的条件。

\begin{resultblock}
\begin{lemma}[具有认证分数区间的分配]
\label{lem:4-assignment-bound}
设 $\widetilde y,\widetilde\mu_\ell$ 为反量化近似，$e=y-\widetilde y$，$f_\ell=\mu_\ell-\widetilde\mu_\ell$。保持质心范数精确时，近似分数
$\widetilde s_\ell=\|\mu_\ell\|^2-2\langle\widetilde y,\widetilde\mu_\ell\rangle$
的实数算术误差由下式界定：
\begin{equation}
 E_\ell=2\bigl(\|y\|\|f_\ell\|+\|e\|\|\widetilde\mu_\ell\|\bigr).
 \label{eq:4-assignment-bound}
\end{equation}
给定包含每个目标分数的有效区间 $[L_\ell,U_\ell]$，仅丢弃 $L_\ell>\min_j U_j$，并对其余分数精确求解，即可保留最小值及规定的并列规则。
\end{lemma}
\begin{proof}
使用 $e=y-\widetilde y$ 和 $f_\ell=\mu_\ell-\widetilde\mu_\ell$，
\[
 \begin{aligned}
 \langle y,\mu_\ell\rangle-
 \langle\widetilde y,\widetilde\mu_\ell\rangle
 &=\langle y,f_\ell\rangle
   +\langle e,\widetilde\mu_\ell\rangle,\\
 |s_\ell-\widetilde s_\ell|
 &\leq2\bigl(\|y\|\|f_\ell\|
       +\|e\|\|\widetilde\mu_\ell\|\bigr)=E_\ell .
 \end{aligned}
\]
对 $j^*\in\arg\min_j U_j$，
\[
 L_\ell>U_{j^*}
 \Longrightarrow
 s_\ell\geq L_\ell>U_{j^*}\geq s_{j^*}.
\]
严格不等号排除 $\ell$；等号情况仍可入选。
\end{proof}
\end{resultblock}

有效分数区间是关键条件：若无累加误差界，FP32 安全余量不能证明分配结果逐位一致。存档 FP16 尾部分配比较在 Vogue 的 0/14,838、BGE 的 6/31,748 和 Stella 的 3/16,359 个分配中报告差异。冷构建日志未标明实际采用的分配精度。

\subsection{残差拟合与迭代预算}

残差相对主码重建定义，即 $r=\Pi x-\widehat p(x)$。残差拟合使用 100,000 行。在每个子空间内，GPU 排序和 Lloyd 更新拟合 $K_r=2^{B_r}$ 个有序标量等级，在各坐标间共享，最多更新 2,000 次。

较早的十二单元 Vogue／OpenAI-3072 扫描（$B_r=4/8$，probe 为 32/128/256）固定主码／IVF 模型。在 25 次 Lloyd 迭代下，以 100K 行替代全部行最多损失 0.0018 召回率（十一个下降、一个持平）。中间样本量的结果并非单调，因此该证据在这一协议内支持有界残差训练。

历史 100K 训练从 25$\to$2,000 次更新增加了 0.0063--0.0076 召回率。尾部 bin 更新和 BGE／Stella 的结果仍在变化：达到上限并不证明收敛，也不违背精确 Lloyd 的下降性质。20,000 次更新时，历史 Stella 全量／100K 召回率为 0.9921/0.9920。

\subsection{训练工作量与列表不均衡}

训练工作量模型分别考虑变换生成、样本旋转、主码本生成和残差提取。以 $n_r$ 表示残差训练行数，变换部分为
\[
 W_{\rm QR}=O(d^3),\qquad W_{\rm rotate}=O(n_t d^2).
\]
直接粗分配每次迭代成本为 $O(n_t n_{\rm list}d)$，其余部分为
\[
 W_{\rm primary}=O(MKD_s),\qquad
 W_{\rm residual\ extraction}=O(n_rd),
\]
随后进行排序与标量更新。这些量表示操作数。同步诊断表明，BGE／Stella 训练约 83\% 的时间在 IVF k-means（约 2.97\,s）。该诊断使用一个 add 缓冲并强制同步，不同于 \cref{sec:7-7} 中正常的重叠测量。列表长度不均影响后续扫描工作量，对应 \cref{sec:5-5}；八次 IVF 迭代、间隔种子和有限抽样各自对不均衡的影响未单独测量。


\section{IVF 分配、ID 与物理内存布局}
\label{sec:4-3}
\subsection{从数据库 ID 到列表顺序的物理位置}

列表排序使分配给同一 IVF 质心的向量连续存储。输入行的标识为 $i\in\{0,\ldots,N-1\}$；粗分配使用计算得到的质心分数，其数学目标为
\begin{equation}
 a(i)=\operatorname*{arg\,min}_{0\leq\ell<n_{\rm list}}
       \|\Pi x_i-\mu_\ell\|_2^2.
 \label{eq:4-assignment}
\end{equation}
实现采用上文说明的数值分配路径。主码和残差编码不减去 $\mu_{a(i)}$；列表归属独立于这些重建，限制后续搜索空间。

构建器将分配结果与输入 ID 一同排序。令 $\rho(p)$ 为排序后物理位置 $p$ 上的输入 ID，$o_\ell$ 为列表 $\ell$ 的起点。完成的数组满足
\begin{equation}
 0=o_0\leq o_1\leq\cdots\leq o_{n_{\rm list}}=N,
 \qquad p\in[o_\ell,o_{\ell+1})\Longrightarrow a(\rho(p))=\ell.
 \label{eq:4-offsets}
\end{equation}
lower-bound 搜索找到各列表的首个位置，因此相同偏移表示空列表：
\[
 o_\ell=o_{\ell+1}\Longleftrightarrow\text{list }\ell\text{ is empty}.
\]
位置在完成的索引中保持固定，重建后可能改变。所有数组使用同一个 $\rho(p)$，ID 通过 $p\mapsto\rho(p)$ 映射回输入行。两遍 add 暂时构造 $\rho^{-1}$，将新编码直接写入目标位置，称为 scatter。主码并列使用物理位置，最终并列使用数据库 ID。

\par\medskip\noindent\begin{minipage}{\linewidth}
\captionsetup{hypcap=false}
\centering
% Source-level schematic: v69 gather/repack and separate allocations.
% All addresses below are byte offsets relative to each allocation's base.
\begin{tikzpicture}[x=1cm,y=1cm,font=\footnotesize\linespread{1}\selectfont,
 cell/.style={draw=ThesisNavy!50,minimum width=.88cm,minimum height=.48cm,inner sep=1pt},
 panel/.style={draw=ThesisNavy!40,rounded corners=2pt,fill=ThesisNavy!2},
 link/.style={-{Stealth[length=4pt]},ThesisNavy!80,semithick}]
\node[anchor=west,font=\small\bfseries] at (0,1.0) {List ranges and physical positions (schematic, $N=8$)};
\node[anchor=west] at (0,.13) {$o=[0,3,3,8]$};
\node[anchor=west] at (0,-.32) {list 1: $[3,3)=\varnothing$};
\foreach \p in {0,...,7} {
 \node[cell,fill=ThesisNavy!5] (pos\p) at (4.4+.88*\p,-.3) {$\p$};
}
\draw[decorate,decoration={brace,amplitude=3pt}] (3.96,.12)--(6.60,.12)
 node[midway,above=4pt] {list 0: $[0,3)$};
\draw[decorate,decoration={brace,amplitude=3pt}] (6.60,.12)--(11.00,.12)
 node[midway,above=4pt] {list 2: $[3,8)$};
\node[anchor=west,align=left] at (11.55,-.05) {Each range indexes\\all four buffers.};
\node[draw=ThesisNavy,fill=ThesisNavy!12,rounded corners=2pt,
 minimum height=.5cm] (chosen) at (1.25,-1.02) {example $p=3$};
\draw[link] (pos3.south)--(7.04,-.72)-|(chosen.north);
\node[anchor=west] at (4.0,-1.02) {Separate allocations; links below denote an index join.};

% ID allocation.
\draw[panel] (1.05,-1.53) rectangle (14.85,-2.57);
\node[anchor=west,align=left] at (1.3,-2.04) {\textbf{IDs}\\$\mathrm{IDs}[p]=\rho(p)$};
\foreach \p/\id in {0/5,1/1,2/7,3/0,4/6,5/3,6/4,7/2} {
 \node[cell,fill=white] at (4.4+.88*\p,-2.04) {$\id$};
}
\node[cell,draw=ThesisNavy,fill=ThesisNavy!17] (id3) at (7.04,-2.04) {$0$};
\node[anchor=west,align=left] at (11.5,-2.04) {int32: $4p$\\$p=3:\ 12$ bytes};

% Primary allocation, word-group-major. The example shows M >= 8.
\draw[panel] (1.05,-2.88) rectangle (14.85,-5.45);
\node[anchor=west,align=left] at (1.3,-3.62) {\textbf{Primary words}\\group-major\\(shown: $M\geq8$)};
\foreach \j/\yy in {0/-3.36,1/-3.96} {
 \foreach \p in {0,...,7} {
  \node[cell,fill=white] at (4.4+.88*\p,\yy) {$\mathsf w_{\j,\p}$};
 }
 \node[cell,draw=ThesisNavy,fill=ThesisNavy!17] (word\j) at (7.04,\yy) {$\mathsf w_{\j,3}$};
}
\node[anchor=west,align=left] at (11.5,-3.65) {$4(jN+p)$\\$p=3,N=8:$\\$j=0:\ 12$ bytes\\$j=1:\ 44$ bytes};
\draw[link] (4.0,-4.43)--(10.95,-4.43)
 node[midway,below=2pt] {fixed $j$: adjacent $p$ give consecutive 32-bit words};
\node[anchor=west] at (1.3,-5.12)
 {One word: $Z[4j,p]$, $Z[4j+1,p]$, $Z[4j+2,p]$, $Z[4j+3,p]$ (one candidate).};

% Residual allocation; each displayed cell represents a whole row.
\draw[panel] (1.05,-5.76) rectangle (14.85,-7.53);
\node[anchor=west,align=left] at (1.3,-6.33) {\textbf{Residual codes}\\row-major (8-bit)};
\foreach \p in {0,...,7} {
 \node[cell,fill=white] at (4.4+.88*\p,-6.33) {row $\p$};
}
\node[cell,draw=ThesisNavy,fill=ThesisNavy!17] (res3) at (7.04,-6.33) {row 3};
\node[anchor=west,align=left] at (11.5,-6.33) {$pb_r+t$\\$p=3:\ 3b_r+t$};
\node[anchor=west] at (1.3,-7.10) {row 3 contains $b_r$ consecutive bytes:};
\node[draw=ThesisNavy!50,fill=white,minimum height=.45cm,inner sep=3pt]
 (expanded) at (9.05,-7.10) {$R[3,0]\quad R[3,1]\quad\cdots\quad R[3,b_r-1]$};
\draw[link] (res3.south)--(expanded.north west);

% Correction allocation.
\draw[panel] (1.05,-7.84) rectangle (14.85,-9.12);
\node[anchor=west,align=left] at (1.3,-8.48) {\textbf{Correction}\\$\mathrm{correction}[p]$\\$=\gamma[p]$};
\foreach \p in {0,...,7} {
 \node[cell,fill=white] at (4.4+.88*\p,-8.36) {$\gamma[\p]$};
}
\node[cell,draw=ThesisNavy,fill=ThesisNavy!17] (corr3) at (7.04,-8.36) {$\gamma[3]$};
\node[anchor=west,align=left] at (11.5,-8.36) {FP32: $4p$\\$p=3:\ 12$ bytes};

% One physical-position key points into four independent buffers.
\draw[ThesisNavy!80,semithick] (chosen.west)--(.35,-1.02)--(.35,-7.69);
\draw[link] (.35,-1.38)-|(id3.north);
\draw[link] (.35,-2.73)-|(word0.north);
\draw[link] (.35,-5.61)-|(res3.north);
\draw[link] (.35,-7.69)-|(corr3.north);
\node[anchor=west,align=left] at (0,-9.65)
 {All offsets are relative to each buffer's base. Refinement fetches survivor rows only;\\
 survivor positions can be scattered. Highlighted entries share $p=3$, not an interleaved allocation.};
\end{tikzpicture}

\captionof{figure}[GPU 索引的物理布局]{GPU 索引的物理布局.}
\label{fig:4-layout}
\end{minipage}\par\medskip

图中示例取 $N=8$，offsets $o=[0,3,3,8]$ 使 list~1 为空。浅色列标出 $p=3$：四个独立数组均以相同物理位置索引，右侧 byte 偏移相对各自数组的起点。固定 word group $j$ 时，相邻候选对应连续的 32 位 word，使扫描相邻候选的线程可以合并访存；八位残差配置下，每行存放同一向量的 $b_r$ 个编码 bytes。

\subsection{精确地址与表示宽度}

主码表示按子空间组织候选。对主码 byte $Z[m,p]$，地址为
\[
 \operatorname{offset}_{\rm byte}(Z[m,p])=mN+p.
\]
最终打包将同一向量的四个子空间 byte 合为一个 word：
\begin{equation}
 \mathsf w[j,p]=\sum_{t=0}^{3}2^{8t}Z[4j+t,p],\qquad
 \operatorname{addr}_{\rm byte}(\mathsf w[j,p])=4(jN+p).
 \label{eq:4-primary-address}
\end{equation}
这里 $j=0,\ldots,M/4-1$。$\text{word-layout build}\Longrightarrow M\bmod4=0$；其余 $M$ 被拒绝，独立的 byte-layout 编译则保留 byte 表示。每个数组占 $NM$ bytes，重打包暂时同时持有两份副本，成本为 $2NM$。一遍构建将已有编码重排（gather）到 vector-major 的 $[N,M]$ 数组，转置并重打包，逐次释放旧副本；两遍构建直接 scatter 到 subspace-major 存储，再进行相同重打包。搜索时的 lane 映射见 \Cref{sec:5-3}。

残差码保持 row-major，每向量占 $b_r=\lceil dB_r/8\rceil$ bytes。八位配置中，行 $p$ 的坐标 $t$ 位于 byte 偏移 $pb_r+t$；四位路径中，它位于 byte $pb_r+\lfloor t/2\rfloor$，偶数坐标在低四位（low nibble），奇数坐标在高四位（high nibble）。每个物理行有四个 byte 的校正项，存储
\begin{equation}
 \gamma[ p ]=2\langle\widehat p(x_{\rho(p)}),
                              \widehat r(x_{\rho(p)})\rangle.
 \label{eq:4-correction}
\end{equation}
gather 或 scatter 时，校正项必须与该行编码一同移动。按列表扫描随后只为保留候选读取残差。精排使用压缩重建；原始 FP32 重排序还需增加 $4Nd$ 个设备 bytes。第 7 章评估其召回率。


\section{常驻存储、工作区与索引生命周期}
\label{sec:4-4}
\subsection{存储字节与分配字节}

常驻行数据只是设备分配的一部分，模型、查询工作区和临时副本另计。主码以 byte 存储时，行与列表偏移的存储为
\begin{equation}
 S_{\rm rows}=N\left(M+\left\lceil\frac{dB_r}{8}\right\rceil+8\right),
 \qquad S_{\rm offsets}=4(n_{\rm list}+1).
 \label{eq:4-row-storage}
\end{equation}
ID 和校正项贡献 $8N$ bytes。当 $M=d/8$、$B_r=8$ 时，仅编码载荷为
\[
 b_{\rm code}=d/8+d=9d/8=1.125d.
\]
这比 8 位标量量化（SQ8）多 12.5\%，因此图存储和向量编码字节数需要一起比较。

\begin{table}[!htbp]
\centering
\caption[每向量存储模型]{每向量字节数。JHQ 包含编码、ID 和校正项；SQ8 仅含向量编码；CAGRA int8 模型包含向量编码及 32 个四字节邻居。不含 offsets、模型参数、工作区及其他运行时分配。这些组件不是完整进程总量。}
\label{tab:4-bytes}
\footnotesize\setlength{\tabcolsep}{4pt}
\begin{tabular}{@{}rrrrrrrr@{}}
\toprule
$d$ & 主码 & 残差 & ID＋校正 & JHQ & SQ8 & \shortstack{CAGRA\\int8} & FP32 \\
\midrule
768 & 96 & 768 & 8 & 872 & 768 & 896 & 3,072 \\
1,024 & 128 & 1,024 & 8 & 1,160 & 1,024 & 1,152 & 4,096 \\
1,536 & 192 & 1,536 & 8 & 1,736 & 1,536 & 1,664 & 6,144 \\
3,072 & 384 & 3,072 & 8 & 3,464 & 3,072 & 3,200 & 12,288 \\
\bottomrule
\end{tabular}
\end{table}

图度数为 $d_g=32$ 时，CAGRA int8 每行使用 $d+4d_g$ bytes。Stella 的 JHQ 行载荷为 19.20\,GiB，接近 int8 CAGRA 报告的 19,656\,MiB 索引分配，但统计边界不同。改变 $\alpha$ 影响残差读取量，不改变已存储的行字节数。

基本模型存储计入旋转矩阵、主码与残差码本，以及质心和其范数：
\[
 S_{\rm model}
 =4d^2+4MKD_s+4MK_r+4n_{\rm list}(d+1)
 \quad\text{bytes}.
\]
派生数组另计。查询工作区包含两个查询矩阵、质心分数、LUT、probes 和保留候选。对预留查询容量 $b_{\rm cap}$，可选残差 LUT 增加 $4b_{\rm cap}dK_r$ bytes，默认直接评分。报告的 128-probe 配置为扫描分组预留约 93.8\,MiB，即使未使用也保留该分配（\cref{eq:5-group-workspace}）；这一大小取决于已分配容量。

输入归调用方所有。对 $n_b$ 个流水缓冲、每缓冲暂存 $a_{\rm in}$ 行，页锁定主机暂存与设备排序数组的存储为
\[
 S_{\rm pinned}=4n_ba_{\rm in}d,\qquad
 S_{\rm order}=S_{\rm inverse}=4N.
\]
逆映射是临时数组。分配临时区在 gather 前释放（\cref{eq:4-staging-rows,eq:4-assignment-scratch}），因此主机暂存、常驻数据和构建峰值具有不同生命周期。

\subsection{缓存复用与失效处理}

缓存键不含代码版本和全量数据哈希，因此代码或数据变化后需使缓存失效。参数扫描复用模型，冷构建重新拟合；仅凭召回率相近不能证明模型相同。缓存命中仍需 QR、预热、模型上传和 add。Stella 的 6.67\,s add 包含初始工作区，不含扫描校准与搜索预热，衡量的是编码调用，而非完整重启时间。

改变 $\alpha$ 无需重训。CUDA Graph 重放会复用查询缓冲，工作区在捕获前扩容。固定返回数量时，改变 batch 大小、保留数、probes、lanes、groups 或 bounds 会触发重新捕获。校准可能绕过捕获，索引拥有的缓冲在释放前保留分配。


\section{构建成本与已测规模}
\label{sec:4-5}
\subsection{有界输入 batch 与两遍构建}

若将全库原始和旋转数组都放在 GPU 上，需要 $4Nd+4Nd=8Nd$ bytes。有界暂存限制这些数组，而压缩索引仍需常驻。对 $n_b$ 个流水缓冲、配置的 add 行数 $a_{\rm add}$ 和每数组预算 $S_{\rm stage}$，暂存行数及原始／旋转设备存储为
\begin{equation}
 a_{\rm in}=\max\!\left\{1,\min\!\left(a_{\rm add},
 \left\lfloor\frac{S_{\rm stage}}{4d}\right\rfloor\right)\right\},
 \qquad S_{\rm float}=8n_ba_{\rm in}d.
 \label{eq:4-staging-rows}
\end{equation}
\[
 4d\leq S_{\rm stage},\quad n_b=3,\ S_{\rm stage}=128\,\mathrm{MiB}
 \Longrightarrow S_{\rm float}\leq2n_bS_{\rm stage}=768\,\mathrm{MiB}.
\]
上界只覆盖这些浮点数组，分配和编码器临时区另计。一遍构建按输入顺序编码，再 gather 为列表顺序；两遍构建先分配和排序，再重读并直接编码到最终位置，以额外遍历换取较少的压缩副本。

\par\medskip\noindent\begin{minipage}{\linewidth}
\captionsetup{hypcap=false}
\centering
\begin{tikzpicture}[font=\footnotesize,
 stage/.style={draw=ThesisNavy!55,fill=ThesisNavy!4,rounded corners=2pt,
 text width=2.5cm,align=center,minimum height=11mm,inner sep=3pt},
 arrow/.style={-{Latex[length=2mm]},draw=ThesisNavy!70}]
\node[stage] (a) at (0,0) {Read + rotate};
\node[stage] (b) at (3.15,0) {Assign + encode\\Input-order arrays};
\node[stage] (c) at (6.3,0) {Gather\\Both copies live};
\node[stage] (d) at (9.45,0) {Repack\\Primary words};
\node[stage] (e) at (0,-1.8) {Assign raw rows\\+ sort};
\node[stage] (f) at (3.15,-1.8) {Reread + rotate\\+ encode};
\node[stage] (g) at (6.3,-1.8) {Scatter directly\\List-order arrays};
\node[stage] (h) at (9.45,-1.8) {Repack\\Primary words};
\foreach \x/\y in {a/b,b/c,c/d,e/f,f/g,g/h}\draw[arrow] (\x)--(\y);
\node[anchor=south west,font=\footnotesize\bfseries] at (-1.3,.7) {One pass};
\node[anchor=south west,font=\footnotesize\bfseries] at (-1.3,-1.1) {Two passes};
\end{tikzpicture}

\captionof{figure}[一遍与两遍构建]{一遍与两遍构建.}
\label{fig:2-add-paths}
\end{minipage}\par\medskip

对 Stella，解析模型给出的行数组占用为：一遍 gather 需要 40.59\,GiB，两遍 repack 需要 21.32\,GiB（\cref{tab:4-phase-storage}）。一遍路径同时保留输入顺序与列表顺序的编码；两遍路径直接写入列表顺序。在首遍分配中，质心通过 $\Pi^\top$ 变回原坐标系，因而无需先旋转数据库行。

\subsection{构建各阶段的显存占用}

实现根据行数组及其他分配的预留量，估计一遍构建所需显存：
\begin{equation}
 H_{\rm single}=N(2b_r+3M+20)+2^{30}\quad\text{bytes}.
 \label{eq:4-peak-heuristic}
\end{equation}
十个存活数组贡献
\[
 3NM+2Nb_r+20N=N(3M+2b_r+20).
\]
1\,GiB 预留量是启发式。两遍构建的最终重打包阶段计入 ID、残差、校正项和两份主码，所需存储为
\begin{equation}
 S_{\rm repack,2}=N(b_r+2M+8).
 \label{eq:4-two-pass-repack}
\end{equation}
整个构建峰值还包括模型、排序、上下文、更早的流水／逆置换临时区和后续查询工作区。\Cref{tab:4-phase-storage} 的阶段值单独计算行数组，并非整个进程峰值。

\begin{table}[!htbp]
\centering
\caption[各构建阶段的行数组占用模型]{行数组占用模型（GiB），$d=1{,}024$、$M=128$、$B_r=8$。BGE 使用一遍，Stella 使用两遍。Stella 的一遍列是反事实容量分析。阶段值不含流水、模型、offsets 和运行时成本。}
\label{tab:4-phase-storage}
\small
\begin{tabular}{@{}lrrr@{}}
\toprule
数据集 & 最终行载荷 & 一遍 Gather & 两遍重打包 \\
\midrule
BGE & 10.90 & 23.05 & 12.11 \\
Stella & 19.20 & 40.59 & 21.32 \\
\bottomrule
\end{tabular}
\end{table}

分配乘积缓冲使用独立容量 $a_{\rm assign}$，由 \texttt{JHQ\_ASSIGN\_BATCH} 或源码默认值选择：至少 8,192 个列表时为 32,768，否则为 8,192。其存储为
\begin{equation}
 S_{\rm dots}=4n_b n_{\rm list}a_{\rm assign}\quad\text{bytes}.
 \label{eq:4-assignment-scratch}
\end{equation}
分配容量独立于当前活跃／暂存行数。在 $d=1536/3072$ 下，两个默认值均分配 3\,GiB。BGE／Stella 的 $3\times4=12$\,GiB 分配临时区在 gather 前释放。对 BGE，12\,GiB 临时区与 23.05\,GiB 行数组 gather 不会同时占用显存。峰值是同一时刻存活分配之和的最大值：
\[
 S_{\rm peak}=\max_t\sum_{A\ {\rm live\ at}\ t}|A|.
\]



Stella 冷构建日志选择两遍：\cref{eq:4-peak-heuristic} 的一遍估计为 41.6\,GiB，同一日志报告尾部阶段可用量为 30.6\,GiB。单独采样的常驻进程总量包含上下文和查询工作区：BGE 为 11.93\,GiB，Stella 为 20.44\,GiB（\cref{sec:4-4}）。补充的 20\,ms 设备采样在 BGE add 开头附近观察到 26.08\,GiB 峰值，在 Stella 重打包时观察到 22.58\,GiB 峰值。较小的两遍重打包量与 Stella 较低的实测峰值一致。这些采样设备总量包含上下文，可能遗漏更短暂的峰值。


\section{索引表示与验证}
\label{sec:4-6}
有符号 32 位的 $n_td$ 在 $2^{31}$ 溢出，使尺度归零；不检查版本的缓存再次载入了该值。分块 FP64 范数修复了错误；仅观察路由成功无法发现它。

\subsection{表示一致性检查}

试验导出覆盖 Vogue、arXiv 和 OpenAI-3072。

码本通过可分性、标量等级和排序检查。FP64 循环检查每个导出行的 $\gamma=2\langle\widehat p,\widehat r\rangle$；最大绝对误差为 $1.28\times10^{-7}$--$2.77\times10^{-7}$。抽样 GPU 返回 top-10 分数的最大相对误差为 $2.38\times10^{-7}$ 至 $4.85\times10^{-5}$，使用包含查询范数常数的 $|s_{\rm GPU}-s_{64}|/s_{64}$。这些检查验证导出的表示，而非从原始输入向量重新编码。

试验还在七个工作负载、896 个查询实例上，将提出的残差界与 FP64 重建距离比较，覆盖 72,800,590 个路由候选实例，报告违规数为零。同一批数据也支持 FP64 复合分数恒等式检查。GPU 返回分数检查只覆盖 top-10 槽位，并跳过不在路由映射中的 ID，因此输出集合成员资格尚未完整检查。

规范化命名空间并移除导出方法后，编码器、主码／残差训练及构建源码与评测实现相符。这种源码对应覆盖共享路径；后续模型身份与大数据集两遍构建另行检查。

\subsection{输入编码与数据放置检查}
\label{sec:4-representation-audit}

补充导出器审计在 Vogue、arXiv、OpenAI-3072 和 Stella 上使用共享构建路径。offsets 单调，端点为 0 和 $N$；四个数据集的 ID 均构成所有数据库行的排列。对存储旋转做 FP64 乘积，$\|\Pi\Pi^{\mathsf T}-I\|$ 的最大逐元素误差为 $3.96$--$4.34\times10^{-7}$。抽样范数检查使用 FP64 旋转，而非 GPU GEMM，因此测量的是存储矩阵的几何性质。\Cref{tab:4-representation-audit} 汇总抽样重编码。

\par\medskip\noindent\begin{minipage}{\linewidth}
\centering\captionsetup{hypcap=false}
\captionof{table}[原始输入的 FP64 重编码审计]{原始输入的 FP64 重编码审计。Flips 计数坐标决策；末列为不一致项到残差决策边界的最大距离。结构检查覆盖全部行，重编码使用抽样。}
\label{tab:4-representation-audit}
\scriptsize\setlength{\tabcolsep}{3pt}
\begin{tabular}{@{}lrrrrr@{}}
\toprule
数据集 & 行数 & 正交误差 & 主码翻转 & 残差翻转 & 最大边界距离 \\
\midrule
Vogue & 200,000 & 4.11e-07 & 14 & 5,409 & 1.12e-07 \\
arXiv & 100,000 & 4.11e-07 & 1 & 2,515 & 7.31e-08 \\
OpenAI-3072 & 20,000 & 3.96e-07 & 5,745 & 1,646,806 & 2.67e-05 \\
Stella & 100,000 & 4.34e-07 & 16 & 3,064 & 1.03e-07 \\
\bottomrule
\end{tabular}

\end{minipage}\par\medskip

在 768/1,024 维下，重编码差异均距边界不超过 $1.12\times10^{-7}$。OpenAI-3072 有 2.68\% 的残差不一致，距边界最高 $2.67\times10^{-5}$，与 TF32 的影响相符。以存储主码为条件，避免重复计算翻转。

按 ID 连接后，三个较小数据库的构建编码完全相同，校正项差不超过 $8.94\times10^{-8}$，Vogue／arXiv／OpenAI 的分配差异为 1/1/40。Stella 仅检查两遍构建。没有导出的质心，因此成员归属尚未验证。

\subsection{验证覆盖范围与局限}
\suppressfloats[t]

\par\medskip\noindent\begin{minipage}{\linewidth}
\captionsetup{hypcap=false}
\centering
\captionof{table}[表示约束与审计覆盖]{试验和补充输入审计后的表示约束。检查分别处理几何、算术和行放置。}
\label{tab:4-contracts}
\footnotesize\setlength{\tabcolsep}{4pt}
\begin{tabularx}{\linewidth}{@{}>{\raggedright\arraybackslash}p{1.8cm}*{2}{>{\raggedright\arraybackslash}X}@{}}
\toprule
边界 & 必要性质 & 证据／剩余检查 \\
\midrule
C0：\newline 几何 & 共同变换及有界浮点等距误差 & 四个数据集检查存储矩阵正交性；区分 TF32 GPU 算术 \\
C1：\newline 码本／校正 & Cartesian 等级、有序残差码本、校正符合解码编码 & 三个数据库完成试验检查；全行校正误差见 \cref{sec:4-6} \\
C1：\newline 输入编码 & 存储码对应预期原始输入行 & 四个样本的 FP64 原始输入重编码；量化边界和 TF32 差异 \\
C1：\newline 放置 & Offsets 覆盖各行；ID 构成排列；所有字段随之移动 & 全行 offsets／ID 检查；三个按 ID 连接的路径比较；质心归属未验证 \\
执行 & 有效索引、同步和分数顺序约束 & 存档查询和 Sanitizer 检查；覆盖有限，详见 \cref{sec:7-4} \\
\bottomrule
\end{tabularx}
\end{minipage}\par\medskip

必要条件为
\[
 M\mid d,\quad D_s\mid B,\quad M\bmod4=0.
\]
要求输入有限、计数有效且哨兵正确。ID 为有符号 32 位；地址乘积使用更宽算术。

T4 验证选择（\cref{sec:7-strict-selection}）；质心归属、空列表和对抗尾部仍未认证。

